\documentclass{article}

\usepackage{xparse}

\usepackage{algorithm}
\usepackage{algorithmic}

\usepackage{amsmath}
\usepackage{amssymb}
\usepackage{tabularray}

\RenewDocumentCommand\emptyset{}{\varnothing}

%%%%% DEFINITIONS %%%%%%%%%%%%%%%%%%%%%%%%%%%%%%%%%%%%%%%%%%%%%%%%%%%%%%%%%%%%%%
\NewDocumentCommand\Handlers{}{\textsf{H}}
\NewDocumentCommand\HandlerIO{}{\ensuremath{\mathcal{H}}}
\NewDocumentCommand\Tuples{}{\ensuremath{T}}
\NewDocumentCommand\Database{}{\textsf{D}}
\NewDocumentCommand\Boot{}{\textsf{boot}}
\NewDocumentCommand\Handler{}{\ensuremath{\langle q, f \rangle}}

\NewDocumentCommand\In{}{\textsf{In}}
\NewDocumentCommand\Out{}{\textsf{Out}}

\DeclareMathOperator{\handlers}{handlers}
%%%%%%%%%%%%%%%%%%%%%%%%%%%%%%%%%%%%%%%%%%%%%%%%%%%%%%%%%%%%%%%%%%%%%%%%%%%%%%%%


\begin{document}

\section{Algorithm}

\begin{tblr}{Q[r,$]l}
  \Database_{i}                   & is the database at timestep $i$. \\
  \Handlers_{i}                   & is the set of handlers. \\
  \Tuples \subseteq \Database_{i} & is a subset of objects in the database. \\
  \Handler{} \in \Handlers_{i}    & is a handler in the set of handlers. \\
  q                               & is a \textit{query}, such that: \\
  \Tuples = \Database(q)          & is the set of objects in \Database{} that match $q$. \\
  f                               & is a function from $\Tuples \rightarrow \Tuples$. \\
  \Boot{}                         & is a user-defined bootstrap handler that matches \bullet. \\
  \bullet                         & is the unit object. \\
\end{tblr}

\vspace{1cm}

We can define the state of the database at step $i$ recursively as:

\begin{displaymath}
  \begin{array}{rcl}
    \langle \Handlers_{0}, \Database_{0} \rangle   & = & \langle \left\{ \Boot{} \right\}, \left\{ \bullet \right\} \rangle \\
    \langle \Handlers_{i+1}, \Database_{i+1} \rangle & =
        & \langle \handlers \Database_{i}, \displaystyle\bigcup_{\Handler \in \Handlers_{i}} f(\Database_{i}(q)) \rangle
  \end{array}
\end{displaymath}

Where $\handlers \Database{}$ returns the subset of handlers in $\Database{}$.

\section{Implementation / Optimization}

This is from a previous try at defining this thing. The algorithm above is, I
think, the minimal way to describe the intended semantics of the system, but it
doesn't necessarily lead to a straightforward, efficient Implementation. It also
doesn't capture side-effectful computation, which we will certainly have: this will
be used to process camera video streams, and render vulkan output.

\dots

We start by adding the following definitions:

\begin{tblr}{Q[r,$]l}
  t \in \Tuples                   & is a set of Tuples.     \\
  \langle \Tuples, \Tuples \rangle \in \HandlerIO_{h}
      & are the I/O pairs of tuples for a handler $h$.     \\
\end{tblr}

If we don't differentiate between aggregating and non-aggregating handlers, the
following algorithm should produce efficient results:

\begin{algorithm}
\caption{Reactor Tick}

\begin{algorithmic}
  \FORALL{$(\Handler = h) \in \Handlers$}
  \STATE $\langle \In', \Out' \rangle = \HandlerIO_{h}$
  \STATE $\In = \Database(q)$
  \IF{$\In = \In'$} \STATE skip to next $h$ \ENDIF
  \STATE $\Database' = \Database \setminus \Out'$
  \STATE $\Out = f(\In)$
  \STATE $\Handlers_{h} := \langle \In, \Out \rangle$
  \STATE $\Database := \Database' \cup \Out$
  \ENDFOR
\end{algorithmic}

\end{algorithm}

\end{document}

